% ============================================================
% ch21.tex —— 第 21 章 加不拢的无穷：调和级数
% ============================================================
\chapter{加不拢的无穷：调和级数}

篇四的引擎是一台"素数密度反应堆"（第 22 章装配），它的燃料是一种特殊的无穷级数。本章研究它——顺便了结第 13 章的悬案：\textbf{每一项都趋近零的级数，一定加得拢吗？}

\begin{kaiwei}
\[
1 + \frac12 + \frac13 + \frac14 + \frac15 + \cdots
\]
每一项都越来越小、趋近 $0$。它加得拢吗？（先押注，再往下读。给你两个提示：不押注的读者，永远学不会押错。）
\end{kaiwei}

\section{直觉押错的一注}

绝大多数人第一反应："能加拢"。理由听起来无懈可击：\textbf{项趋近 $0$}，加到后来"几乎没加什么"，总和应该稳定在某个数附近。连数学史上都有名家栽过——十四世纪的计算者奥雷姆之前，普遍相信它收敛。

看数据，摆证据：

\begin{center}
\begin{tabular}{ccccccc}
\toprule
加到 & $10$ 项 & $100$ 项 & $10^3$ 项 & $10^6$ 项 & $10^{10}$ 项 & $10^{43}$ 项 \\
\midrule
部分和 $H_n$ & $2.929$ & $5.187$ & $7.485$ & $14.393$ & $23.603$ & $\approx 100$ \\
\bottomrule
\end{tabular}
\end{center}

爬得极慢——一百万项才爬到 $14.4$；从 $10^6$ 到 $10^{10}$ 项（多加一万亿项），只涨 $9$——但请盯住趋势：\textbf{它从不减速到停止}。部分和序列 $H_n$ 单调上升，且总能再涨一点。\textbf{它是收敛还是发散？}数据本身不裁决（"涨得慢"与"涨到头"肉眼难分），需要证明。

\begin{dingli}[奥雷姆，约 1350 年]
调和级数发散：
\[
1 + \frac12 + \frac13 + \frac14 + \cdots = \infty .
\]
\end{dingli}

\begin{proof}[分块放大术（六百年前的两行刀法）]
把无穷多项按 $2$ 的幂切成一段段"块"，给每块一个\textbf{下限}：
\[
\underbrace{1}_{\text{第 }0\text{ 块}} + \underbrace{\frac12}_{\text{第 }1\text{ 块}} + \underbrace{\frac13 + \frac14}_{\text{第 }2\text{ 块}} + \underbrace{\frac15 + \frac16 + \frac17 + \frac18}_{\text{第 }3\text{ 块}} + \underbrace{\frac19 + \cdots + \frac1{16}}_{\text{第 }4\text{ 块}} + \cdots
\]
逐块验算下限：
\begin{itemize}
\item 第 $2$ 块（两项）：两项都 $\ge \tfrac14$，故块和 $\ge \tfrac14\times2 = \tfrac12$；
\item 第 $3$ 块（四项）：每项 $\ge \tfrac18$，块和 $\ge \tfrac18\times4 = \tfrac12$；
\item 第 $4$ 块（八项）：每项 $\ge \tfrac1{16}$，块和 $\ge \tfrac1{16}\times8 = \tfrac12$；
\item 一般地，第 $k$ 块（从 $\tfrac{1}{2^{k-1}+1}$ 到 $\tfrac{1}{2^k}$，共 $2^{k-1}$ 项）：每项 $\ge \tfrac{1}{2^k}$，块和 $\ge \dfrac{1}{2^k}\times 2^{k-1} = \dfrac12$。
\end{itemize}
\textbf{无穷多块，每块至少贡献 $\tfrac12$}：
\[
H_{2^k} \;\ge\; \underbrace{\frac12 + \frac12 + \cdots + \frac12}_{k \text{ 块}} = \frac{k}{2} \;\xrightarrow{k\to\infty}\; \infty .
\]
部分和无界，而正项级数部分和单调上升——单调上升无上界，发散。证毕。
\end{proof}

干净得像刀。注意它同时拆穿了直觉的病根：$\tfrac1{1000000}$ 确实近乎为零，\textbf{但想凑出这么多项，前面要垫几十万个"不算太小"的项}——\textbf{项在变小，项的个数在以同样节奏变多}，两股力量打平，且永远打平。收敛需要"变小"跑赢"变多"，调和级数是平局的标本。

\section{精确丈量：$H_n$ 长什么样}

发散，但"以什么速度"发散？这决定它能不能当尺子用。三组对账：

\begin{center}
\begin{tabular}{cccc}
\toprule
$n$ & $H_n$ & $\ln n$ & $H_n - \ln n$ \\
\midrule
$10$ & $2.9290$ & $2.3026$ & $0.6264$ \\
$100$ & $5.1874$ & $4.6052$ & $0.5822$ \\
$10^6$ & $14.3927$ & $13.8155$ & $0.5772$ \\
$10^{10}$ & $23.6027$ & $23.0259$ & $0.5772$ \\
\bottomrule
\end{tabular}
\end{center}

第三列暴露了天机：$H_n - \ln n$ 落在 $0.5772$ 附近\textbf{纹丝不动}（越来越稳）。结论：
\[
H_n \approx \ln n + \gamma, \qquad \gamma \approx 0.57721566\cdots
\]
$\gamma$（伽马）叫\textbf{欧拉常数}——一个与 $\mathrm{e}$、$\pi$ 同级别的常数，但身世成谜：\textbf{人类至今不知道它是不是无理数}（三百年的悬案，比黎曼猜想还早）。对我们而言，要紧的是这条结论的翻译：
\[
\boxed{\;\text{调和级数爬行的速度，恰好是对数的速度}\;}
\]
加到超过 $100$ 需要 $\mathrm{e}^{100 - \gamma} \approx 1.5\times10^{43}$ 项——$10^{43}$ 已超过可观测宇宙恒星数（$\sim10^{24}$）十几个数量级；想加到 $200$，需要约 $10^{87}$ 项，连可观测宇宙的原子数（$\sim10^{80}$）都不够给每一项发一个住址。换一种说法：\textbf{每多爬 $1$，需要把项数乘 $\mathrm{e}\approx2.7$ 倍}——指数级的耐心，换对数级的回报。

\section{对照组：分母的性格决定命运}

把分母换几种性格，命运的分化立竿见影：

\begin{center}
\begin{tabular}{cccc}
\toprule
级数 & 分母性格 & 结局 & 凭据 \\
\midrule
$\sum\dfrac1n$ & $+1$ 爬坡 & 发散（$\ln$ 速） & 奥雷姆分块 \\
$\sum\dfrac{1}{n\ln n}$ & 爬坡再慢一档 & \textbf{仍发散}（$\ln\ln$ 速） & 分块术改版（烧脑时刻） \\
$\sum\dfrac{1}{n(\ln n)^2}$ & 再慢两档 & 收敛 & 比较判别 \\
$\sum\dfrac{1}{n^2}$ & 平方加速 & 收敛于 $\dfrac{\pi^2}{6}$ & 巴塞尔问题（第 24 章） \\
$\sum\dfrac1{n!}$ & 阶乘爆炸 & 收敛于 $\mathrm{e}-1$（飞速） & 第 19 章 \\
\bottomrule
\end{tabular}
\end{center}

第二行值得单独看一眼：$\sum\tfrac{1}{n\ln n}$ 的分母比 $\ln n$ 还大，\textbf{居然还发散}——"平局"的边缘如此靠近收敛，靠 $\ln\ln$（对数套对数）才勉强打破。这张表的分界线（$p$ 次幂：$\sum \tfrac1{n^p}$ 当 $p > 1$ 收敛、$p\le1$ 发散）是大学"级数判别法"的核心结论，其哲学在这里已经全裸呈现：\textbf{收敛与否不取决于"项是否变小"，取决于变小的速度是否跑赢计数}。

\begin{cidian}
本章主角 $\sum\tfrac1n$ 叫\textbf{调和级数}（harmonic series）——名字来自音乐：一根弦的 $\tfrac12, \tfrac13, \tfrac14,\dots$ 处分出的泛音，波长恰成调和比例 $1:\tfrac12:\tfrac13:\cdots$（毕达哥拉斯学派发现弦长比 $\tfrac23$、$\tfrac34$ 给出和谐音程——"调和"由乐理进入数学）。部分和 $H_n$ 叫\textbf{调和数}。
\end{cidian}

\section{为什么本书需要发散}

回到素数。发问：\textbf{把调和级数里的合数全部删掉，只留素数项}：
\[
\frac12 + \frac13 + \frac15 + \frac17 + \frac1{11} + \frac1{13} + \cdots \;=\; ?
\]
素数比整数稀那么多（占比 $2\%\to0$），删剩的级数是"稠得能加到无穷"，还是"稀得加得拢"？这个问题是素数密度的\textbf{终极化验}：

- 若\textbf{收敛}：素数稀到"质量有限"，在整数世界可忽略——素数定理的图景将完全不同；
- 若\textbf{发散}：素数虽稀，但"质量无穷"，撑得起一片天。

第 22 章欧拉将证明：\textbf{发散}——而且发散的方式（$\ln\ln x$ 速，慢到令人发指）与素数密度 $\tfrac1{\ln x}$ 严丝合缝。\textbf{一个发散的级数，即将变成丈量素数的尺子}。本章的调和级数是那把尺子的刻度原型：$H_n \sim \ln n$，素数倒数和 $\sim \ln\ln x$——\textbf{尺子的尺子}，层层对数，这就是解析数论的地形。

\begin{dongshou}
\begin{itemize}
  \yisi 手工执行分块术到第 $5$ 块：计算 $\tfrac13+\tfrac14 \ge \tfrac12$、$\tfrac15+\cdots+\tfrac18 \ge \tfrac12$、$\tfrac19+\cdots+\tfrac1{16} \ge \tfrac12$ 的每一步（列出每项与下限 $\tfrac1{8}$、$\tfrac1{16}$ 的比较），并算出 $H_{16} \ge \tfrac52 = 2.5$（真值 $3.3807$，下限有效但宽——体会"下限证明"只保方向不保精度）。
  \yier 证明"删掉所有含数字 $9$ 的项"之后，级数 $\sum'\tfrac1n$ \textbf{收敛}。（提示：$1$--$10^k$ 里不含 $9$ 的数有 $8^k$ 个（每位 $8$ 种选择），其中最大的一个是 $\tfrac{1}{10^{k-1}}$ 量级。分块：第 $k$ 位段（$10^{k-1}$ 到 $10^k$）的删后总和 $\le 8^k\times\tfrac1{10^{k-1}} = 10\cdot(0.8)^k$——几何级数收敛！对比：只删九的倍数照样发散，删"含九的数"却收敛——\textbf{按数字长相删比按整除性删狠得多}。）
  \yier 交错调和级数 $1 - \tfrac12 + \tfrac13 - \tfrac14 + \cdots$ 收敛在 $0.5$ 与 $1$ 之间：把相邻两项括号成 $(1-\tfrac12) + (\tfrac13-\tfrac14)+\cdots$，每对为正；再用"部分和上下夹逼"（奇数项部分和递减、偶数项递增、彼此靠近）说清收敛。它的精确值 $\ln 2 \approx 0.6931$——第 19 章烧脑时刻你算过前八项 $0.6345$、前十六项 $0.6629$，现在你知道了它奔向谁。
\end{itemize}
\end{dongshou}

\begin{shaonao}
\textbf{分块术改版：证明 $\sum_{n=2}^{\infty}\dfrac{1}{n\ln n}$ 发散。}工具只需要一点"离散微积分"：

(1) 先建立下限：对每个 $n \ge 2$，$\ln(n+1) - \ln n = \ln\left(1+\tfrac1n\right) \le \tfrac1n$（提示：$\ln(1+u) \le u$，来自 $\mathrm{e}^u \ge 1+u$，两边取 $\ln$）。

(2) 于是对区间 $[n, n+1)$ 上的积分 $\displaystyle\int_n^{n+1}\frac{dt}{t} = \ln(n+1)-\ln n \le \frac1n$——\textbf{矩形面积 $\ge$ 曲线下面积}（高取左端点，函数递减时矩形偏高）。

(3) 把第 (2) 步用到分母 $n\ln n$ 上：块 $[2^k, 2^{k+1})$ 内，$\tfrac{1}{t\ln t}$ 递减，故
\[
\sum_{n=2^k}^{2^{k+1}-1}\frac{1}{n\ln n} \;\ge\; \int_{2^k}^{2^{k+1}}\frac{dt}{t\ln t} = \Big[\ln(\ln t)\Big]_{2^k}^{2^{k+1}} = \ln(k+1) - \ln k .
\]
（换元 $u = \ln t$，$\int \tfrac{dt}{t\ln t} = \ln\ln t$——把 $\ln$ 当新变量，幂法则倒抄。）

(4) 把 $k = 1, 2, 3, \dots$ 的块下限相加：$\sum_k [\ln(k+1)-\ln k] = \ln(N) \to \infty$（望远镜求和，相邻项相消）。发散，且发散速度正是 $\ln\ln$。

做完这四步，你已亲手造出第 22 章"素数密度反应堆"的核心零件：\textbf{用积分给级数定速}的技术（积分比较判别法），欧拉当年碾过素数倒数和的履带，就是你刚发明的这块。
\end{shaonao}

\begin{chengshi}
"$H_n - \ln n$ 收敛到 $\gamma$"的严格证明（单调有界定理：$H_n - \ln n$ 递减有下界）只需半页，我们直接引用结论；$\gamma$ 的无理性是未解之谜（三百年来最好的结果：若 $\gamma$ 是有理数 $\tfrac pq$，则 $q > 10^{242080}$）。"$p$-级数判别"（$p > 1$ 收敛）的证明用积分比较（烧脑时刻已发预览）。另外注意"部分和单调上升且无上界 $\Rightarrow$ 发散"只对\textbf{正项}级数成立——交错级数的部分和会自己刹车，动手实验第三题就是活例。
\end{chengshi}

燃料备齐，履带已造。下一章，欧拉点火：把"加法的世界"和"乘法的世界"焊在一起，看素数从齿轮里现出原形。
